% Source coverage: physical PDF pp. 439--443 (printed pp. 416--420),
% beginning at the Problems heading and continuing through the complete
% Chapter 22 reference list.
% Problems: 4. References: 37 entries (1--26 plus 6a, 7a, 8a, 10a, 10b,
% 13a, 16a, 18a, 18b, 18c, and 19a). Unnumbered displays, footnotes,
% figures, tables, and centered dividers: none.
% Visible source anomalies retained verbatim include missing spaces in
% Refs. 10a and 15, the double comma after "E. Witten" and the missing
% comma before "Phys. Lett." in Ref. 16, "Nuc. Phys." in Ref. 15, and
% the anomalous volume "96" in Ref. 11.

\chapterbackmatter{Problems}

\begin{enumerate}
\item Calculate the rate of the decay process $\eta\to\gamma+\gamma$, to
leading order in $m_s$, with $m_u=m_d=0$.

\item Consider a chiral $SU(3)$ symmetry under which the left-handed parts
of the spin $\frac{1}{2}$ fields of a fermion-number-conserving theory form
$N$ defining representations $\boldsymbol{3}$ of $SU(3)$, while the
right-handed parts are all singlets. Evaluate the anomaly in the $SU(3)$
symmetry. What is the anomaly if we add $M$ fermion fields whose left-handed
parts are singlets, and whose right-handed parts transform as symmetric
traceless second-rank $SU(3)$ tensors?

\item Find a solution of the 't Hooft anomaly matching conditions
\hyperref[eq:22.5.5]{(22.5.5)} and \hyperref[eq:22.5.6]{(22.5.6)} for the
case of $n=4$ flavors. Find a solution for $n=2$ flavors other than the one
given in the text.

\item Derive the Zinn-Justin equation from the quantum master equation,
without assuming that $\Delta S=0$.
\end{enumerate}

\chapterbackmatter{References}

\begin{enumerate}
\item \phantomsection\label{ch22-ref-1}
J. Steinberger, \emph{Phys. Rev.} \textbf{76}, 1180 (1949); Also see
R. J. Finkelstein, \emph{Phys. Rev.} \textbf{72}, 415 (1949);
H. Fukuda and Y. Miyamoto, \emph{Prog. Theor. Phys.} \textbf{4}, 347
(1949); J. Schwinger, \emph{Phys. Rev.} \textbf{82}, 664 (1951);
L. Rosenberg, \emph{Phys. Rev.} \textbf{129}, 2786 (1963). A rough
estimate of this decay rate was made by S. Sakata and Y. Tanikawa,
\emph{Phys. Rev.} \textbf{57}, 548 (1940), before the experimental
discovery of the $\pi^0$ meson.

\item \phantomsection\label{ch22-ref-2}
D. G. Sutherland, \emph{Nucl. Phys.} \textbf{B2}, 433 (1967)

\item \phantomsection\label{ch22-ref-3}
M. Veltman, \emph{Proc. Roy. Soc.} \textbf{A301}, 107 (1967).

\item \phantomsection\label{ch22-ref-4}
J. S. Bell and R. Jackiw, \emph{Nuovo Cimento} \textbf{60A}, 47 (1969).
Also see R. Jackiw, in \emph{Lectures on Current Algebra and its
Applications}, (Princeton Press, Princeton, 1972).

\item \phantomsection\label{ch22-ref-5}
S. Adler, \emph{Phys. Rev.} \textbf{177}, 2426 (1969). Also see S. Adler,
in \emph{Lectures on Elementary Particles and Quantum Field Theory --
1970 Brandeis University Summer Institute in Theoretical Physics}, eds.
S. Deser, M. Grisaru, and H. Pendleton (MIT Press, Cambridge, MA, 1970):
Volume I.

\item \phantomsection\label{ch22-ref-6}
K. Fujikawa, \emph{Phys. Rev. Lett.} \textbf{42}, 1195 (1979).

\item[6a.] \phantomsection\label{ch22-ref-6a}
M. F. Atiyah and I. M. Singer, \emph{Proc. Nat. Acad. Sci.} \textbf{81},
2597 (1984).

\item \phantomsection\label{ch22-ref-7}
L. Alvarez-Gaum\'e and P. Ginsparg, \emph{Ann. Phys.} \textbf{161}, 423
(1985).

\item[7a.] \phantomsection\label{ch22-ref-7a}
For a more detailed argument, see Y. Frishman, A. Schwimmer, T. Banks,
and S. Yankielowicz, \emph{Nucl. Phys.} \textbf{B177}, 157 (1981).

\item \phantomsection\label{ch22-ref-8}
W. A. Bardeen, \emph{Phys. Rev.} \textbf{184}, 1848 (1969).

\item[8a.] \phantomsection\label{ch22-ref-8a}
See, e.g., B. Zumino, Y-S. Wu, and A. Zee, \emph{Nucl. Phys.}
\textbf{B239}, 477 (1984).

\item \phantomsection\label{ch22-ref-9}
S. L. Adler and W. A. Bardeen, \emph{Phys. Rev.} \textbf{182}, 1517
(1969).

\item \phantomsection\label{ch22-ref-10}
See, e.g., H. Georgi,\emph{Lie Algebras in Particle Physics}
(Benjamin--Cummings, Reading, MA, 1982): pp. 15, 198. The original
reference is I. Schur, \emph{Sitz. Preuss. Akad.}, p. 406 (1905).

\item[10a.] \phantomsection\label{ch22-ref-10a}
M. F. Atiyah and I. M. Singer, \hyperref[ch22-ref-6a]{Ref.~6a};
B. Zumino, in \emph{Relativity, Groups, and Topology II}, eds.
B. S. De Witt and R. Stora (North-Holland, Amsterdam, 1984);
L. Alvarez-Gaum\'e and P. Ginsparg, \hyperref[ch22-ref-7]{Ref.~7};
W. Bardeen and B. Zumino,\emph{Nucl. Phys.} \textbf{B244}, 421 (1984);
R. Stora, in \emph{Progress in Gauge Field Theory}, eds. G. 't Hooft
\emph{et al.}(Plenum, New York, 1984): 543; O. Alvarez, I. Singer,and
B. Zumino, \emph{Comm. Math. Phys} \textbf{96}, 409 (1984);
L. Alvarez-Gaum\'e and P. Ginsparg, \emph{Nucl. Phys.} \textbf{B262},
439 (1985); J. Manes, R. Stora, and B. Zumino,
\emph{Commun. Math. Phys.} \textbf{102}, 157 (1985); B. Zumino,
\emph{Nucl. Phys.} \textbf{B253}, 477 (1985); J. M. Bismut and
D. S. Freed, \emph{Commun. Math. Phys.} \textbf{106}, 159 (1986);
\textbf{107}, 103 (1986); D. S. Freed, \emph{Commun. Math. Phys.}
\textbf{107}, 483 (1986).

\item[10b.] \phantomsection\label{ch22-ref-10b}
E. Witten, \emph{Phys. Lett.} \textbf{117 B}, 324 (1982).

\item \phantomsection\label{ch22-ref-11}
D. J. Gross and R. Jackiw, \emph{Phys. Rev.} \textbf{96}, 477 (1969).

\item \phantomsection\label{ch22-ref-12}
H. Georgi and S. L. Glashow, \emph{Phys. Rev.} \textbf{D6}, 429 (1972).

\item \phantomsection\label{ch22-ref-13}
M. L. Mehta, \emph{J. Math. Phys.} \textbf{7}, 1824 (1966);
M. L. Mehta and P. K. Srivastava, \emph{J. Math. Phys.} \textbf{7},
1833 (1966).

\item[13a.] \phantomsection\label{ch22-ref-13a}
The cancellation of anomalies in the four-quark version of the standard
electroweak theory was shown by C. Bouchiat, J. Iliopoulos, and Ph. Meyer,
\emph{Phys. Lett} \textbf{38B}, 519 (1972); S. Weinberg, in
\emph{Fundamental Interactions in Physics and Astrophysics}, eds.
G. Iverson \emph{et al.} (Plenum Press, New York, 1973): p. 157.

\item \phantomsection\label{ch22-ref-14}
H. Georgi, in \emph{Particles and Fields -- 1974}, ed. C. Carlson
(Amer. Inst. of Physics, New York, 1975).

\item \phantomsection\label{ch22-ref-15}
R. Delbourgo and A. Salam, \emph{Phys. Lett.} \textbf{40B}, 381 (1972);
T. Eguchi and P. Freund, \emph{Phys. Rev. Lett.} \textbf{37}, 1251
(1976). Also see N. K. Nielsen, M. T. Grisaru, R. Romer, and
P. van Nieuwenhuizen, \emph{Nucl. Phys.} \textbf{B140}, 477 (1978);
M. J. Perry, \emph{Nucl. Phys.} \textbf{B143}, 114 (1978);
S. W. Hawking and C. Pope, \emph{Nucl. Phys.} \textbf{B146}, 381 (1978);
S. M. Christensen and M. J. Duff, \emph{Phys. Lett.} \textbf{76B}, 571
(1978); R. Critchley, \emph{Phys. Lett.} \textbf{78B}, 410 (1978);
A. J. Hanson and H. Romer, \emph{Phys. Lett.} \textbf{80B}, 58 (1978).
For a general review,see T. Eguchi, P. B. Gilkey, and A. J. Hanson,
\emph{Phys. Rep.} \textbf{66}, 213 (1980). In $4n+2$ dimensions there is
also an anomaly in the divergence of the energy momentum tensor; see
L. Alvarez-Gaum\'e and E. Witten, \emph{Nuc. Phys.} \textbf{B234}, 269
(1984).

\item \phantomsection\label{ch22-ref-16}
G. 't Hooft, lecture given at the Cargese Summer Institute, 1979, in
\emph{Recent Developments in Gauge Theories}, eds. G. 't Hooft
\emph{et al.} (Plenum, New York, 1980), reprinted in
\emph{Dynamical Gauge Symmetry Breaking}, eds. E. Farhi and R. Jackiw
(World Scientific, Singapore, 1982), and in G. 't Hooft,
\emph{Under the Spell of the Gauge Principle} (World Scientific,
Singapore, 1994). Also see S. Dimopoulos, S. Raby, and L. Susskind,
\emph{Nucl. Phys.} \textbf{B173}, 208 (1980); S. Coleman and
E. Witten,, \emph{Phys. Rev. Lett.} \textbf{45}, 1000 (1980);
Y. Frishman, A. Schwimmer, T. Banks, and S. Yankielowicz,
\emph{Nucl. Phys.} \textbf{B177}, 157 (1981); A. Zee, Univ. of
Pennsylvania report, 1980 (unpublished); R. Barbieri, L. Maiani, and
R. Petronzio, \emph{Phys. Lett.} \textbf{96B}, 63 (1980); G. Farrar,
\emph{Phys. Lett.} \textbf{96B}, 273 (1980); R. Chanda and P. Roy
\emph{Phys. Lett.} \textbf{99B}, 453 (1981).

\item[16a.] \phantomsection\label{ch22-ref-16a}
S. Weinberg and E. Witten, \emph{Phys. Lett.} \textbf{96B}, 59 (1980).

\item \phantomsection\label{ch22-ref-17}
J. Preskill and S. Weinberg, \emph{Phys. Rev.} \textbf{D24}, 1059 (1981).

\item \phantomsection\label{ch22-ref-18}
J. Wess and B. Zumino, \emph{Phys. Lett.} \textbf{37B}, 95 (1971).

\item[18a.] \phantomsection\label{ch22-ref-18a}
R. Stora, in \emph{Progress in Gauge Field Theory}, eds. G. 't Hooft
\emph{et al.} (Plenum, New York, 1984):543; B. Zumino, in
\emph{Relativity, Groups and Topology II}, eds. B. S. De Witt and
R. Stora (Elsevier, Amsterdam, 1984):1293; J. Ma\~nes, R. Stora, and
B. Zumino, \emph{Commun. Math. Phys.} \textbf{102}, 157 (1985).

\item[18b.] \phantomsection\label{ch22-ref-18b}
J. Schwinger, \emph{Phys. Rev. Lett.} \textbf{3}, 296 (1959).

\item[18c.] \phantomsection\label{ch22-ref-18c}
B. Zumino, \emph{Nucl. Phys.} \textbf{B253}, 477 (1985).

\item \phantomsection\label{ch22-ref-19}
J. A. Dixon, unpublished preprints (1976--9); \emph{Commun. Math. Phys.}
\textbf{139}, 495 (1991); F. Brandt, N. Dragon, and M. Kreuzer,
\emph{Nucl. Phys.} \textbf{B332}, 224 (1990); M. Dubois-Violette,
M. Henneaux, M. Talon, and C. M. Viallet, \emph{Phys. Lett.}
\textbf{B289}, 361 (1992); M. Dubois-Violette, M. Talon, and
C. M. Viallet, \emph{Phys. Lett.} \textbf{B158}, 231 (1985);
\emph{Commun. Math. Phys.} \textbf{102}, 105 (1985); J. Ma\~nes and
B. Zumino, in \emph{Supersymmetry and its Applications: Superstrings,
Anomalies, and Supergravity}, eds. G. W. Gibbons, S. W. Hawking, and
P. K. Townsend (Cambridge University Press, Cambridge, 1986).

\item[19a.] \phantomsection\label{ch22-ref-19a}
See, e.g., H. Georgi, \hyperref[ch22-ref-10]{Ref.~10}: Eq. (XXVII.9).

\item \phantomsection\label{ch22-ref-20}
R. Stora, in \emph{New Directions in Quantum Field Theory and Statistical
Mechanics -- Lectures at the 1976 Cargese Summer School}, eds. M. L\'evy
and P. Mitter (Plenum, New York, 1977).

\item \phantomsection\label{ch22-ref-21}
J. A. Dixon, \hyperref[ch22-ref-19]{Ref.~19}.

\item \phantomsection\label{ch22-ref-22}
G. Barnich and M. Henneaux, \emph{Phys. Rev. Lett.} \textbf{72}, 1588
(1994); G. Barnich, F. Brandt, and M. Henneaux,
\emph{Commun. Math. Phys.} \textbf{174}, 57, 93 (1995).

\item \phantomsection\label{ch22-ref-23}
W. Troost, P. van Nieuwenhuizen, and A. Van Proeyen,
\emph{Nucl. Phys.} \textbf{B333}, 727 (1990).

\item \phantomsection\label{ch22-ref-24}
I learned about this proof from B. Zumino, private communication.

\item \phantomsection\label{ch22-ref-25}
E. Witten, \emph{Nucl. Phys.} \textbf{B223}, 422 (1983).

\item \phantomsection\label{ch22-ref-26}
C-S. Chu, P-M. Ho, and B. Zumino, Berkeley preprint LBL-38276,
UCB-PTH-96/05, hep-th/9602093, to be published (1996).
\end{enumerate}
